\documentclass[twoside,12pt]{book}

\usepackage[utf8]{inputenc}
\usepackage[T1]{fontenc}
\usepackage{mathpazo}
\usepackage{amsmath,amssymb,amsthm}
\usepackage{graphicx}
\usepackage{listings}
\usepackage{xcolor}
\usepackage{enumitem}
\usepackage{tikz}
\usepackage{algorithm}
\usepackage{algpseudocode}
\usepackage{booktabs}
\usepackage{geometry}
\usepackage{fancyhdr}
\usepackage{titlesec}
\usepackage[colorlinks=true, linkcolor=blue, citecolor=blue, urlcolor=blue, plainpages=false, pdfpagelabels]{hyperref}

\geometry{
  papersize={6in,9in},
  margin=0.75in,
  top=0.75in,bottom=0.75in,
  inner=0.85in,outer=0.65in,
  headheight=14pt
}

\pagestyle{fancy}
\fancyhf{}
\fancyhead[LE]{\small\textit{7. Markov Chains and Random Walks}}
\fancyhead[RO]{\small\textit{Randomized Algorithms}}
\fancyfoot[C]{\thepage}
\renewcommand{\headrulewidth}{0.4pt}

\definecolor{codegray}{rgb}{0.45,0.45,0.45}
\lstdefinestyle{cppstyle}{
  basicstyle=\ttfamily\scriptsize,
  keywordstyle=\bfseries,
  commentstyle=\itshape\color{codegray},
  numbers=left,
  numberstyle=\tiny\color{codegray},
  numbersep=8pt,
  frame=single,
  framerule=0.4pt,
  rulecolor=\color{codegray},
  backgroundcolor=\color{gray!5},
  breaklines=true,
  tabsize=4,
  showstringspaces=false,
  language=C++,
  morekeywords={nullptr,size_t,auto,const,constexpr,
    unordered_map,vector,pair,mt19937,
    uniform_int_distribution,INT_MAX}
}
\lstset{style=cppstyle}

\theoremstyle{plain}
\newtheorem{theorem}{Theorem}[chapter]
\newtheorem{lemma}[theorem]{Lemma}
\newtheorem{corollary}[theorem]{Corollary}
\newtheorem{proposition}[theorem]{Proposition}
\theoremstyle{definition}
\newtheorem{definition}[theorem]{Definition}
\newtheorem{example}[theorem]{Example}
\newtheorem{remark}[theorem]{Remark}
\theoremstyle{remark}
\newtheorem{exercise}[theorem]{Exercise}
\newtheorem{problem}[theorem]{Problem}
\newenvironment{cpplisting}{\begin{center}\small}{\end{center}}

\newcommand{\Exp}{\operatorname{\mathbb{E}}}
\newcommand{\Var}{\operatorname{Var}}
\newcommand{\Cov}{\operatorname{Cov}}

\begin{document}

\chapter*{7 Markov Chains and Random Walks}
\addcontentsline{toc}{chapter}{7 Markov Chains and Random Walks}
\markboth{7 Markov Chains and Random Walks}{Randomized Algorithms}
\label{ch:markov-chains}

Markov chains and random walks are among the most versatile tools in the
theory of randomized algorithms. A Markov chain describes a process that
moves among states of a system according to a memoryless transition rule.
When the chain is ergodic---irreducible and aperiodic---it possesses a
unique stationary distribution to which it converges regardless of its
starting state. This convergence property underlies a powerful paradigm:
to sample from a complex distribution, construct a Markov chain whose
stationary distribution is the target, run the chain long enough to
``mix,'' and then read off a sample.

Random walks on graphs are the canonical special case. Given a graph, the
walker at each step moves to a uniformly random neighbor. The resulting
chain encodes deep structural information about the graph: connectivity,
bottlenecks, and expansion. The interplay between random walks and
electrical networks connects probability to combinatorial optimization.
Applications range from testing graph connectivity and estimating cover
times, to sampling almost-uniformly from the solutions of combinatorial
problems via rapidly mixing chains on expander graphs.

In this chapter, we begin with a concrete motivating example---solving
2-SAT by random walk---and then develop the general theory of Markov
chains, random walks on graphs, electrical network analogies, cover
times, expander-based sampling, and probability amplification.

%======================================================================
\section{A 2-SAT Example}
\label{sec:2sat-example}
%======================================================================

The 2-satisfiability problem asks: given a Boolean formula $\phi$ in
conjunctive normal form where each clause has exactly two literals, does
there exist a satisfying assignment? While 2-SAT is solvable in
deterministic polynomial time by a clever graph algorithm, the random
walk approach to 2-SAT provides an elegant introduction to the use of
Markov chains for combinatorial search and foreshadows the general
paradigm of random-walk-based satisfiability algorithms.

\subsection{The Walksat Algorithm for 2-SAT}

Consider a 2-CNF formula $\phi$ with $n$ variables $x_1, \ldots, x_n$
and $m$ clauses. The Walksat algorithm proceeds as follows: begin with an
arbitrary assignment, and repeat. If the current assignment satisfies
$\phi$, halt. Otherwise, pick an unsatisfied clause uniformly at random,
and then flip the value of one of its two variables, chosen uniformly at
random.

To analyze this algorithm, we consider the number of satisfied clauses.
Let $S$ denote the current count of satisfied clauses. Each unsatisfied
clause $(l_1 \vee l_2)$ has both $l_1$ and $l_2$ false. Flipping either
literal makes the clause satisfied, increasing $S$ by~$1$. The other
variable might appear in other clauses, so flipping it could decrease $S$
by some amount. However, the expected change in $S$ can be bounded
carefully.

\begin{theorem}[2-SAT Random Walk]\label{thm:2sat-walk}
If the formula $\phi$ is satisfiable, then Walksat finds a satisfying
assignment in expected $O(n^2)$ steps.
\end{theorem}

\begin{proof}
Let $x^*$ be a fixed satisfying assignment. Define the distance
$d(t) = |\{i : x_i(t) \neq x_i^*\}|$ as the number of variables on which
the current assignment disagrees with $x^*$. Initially $d(0) \leq n$.
When an unsatisfied clause $(l_1 \vee l_2)$ is chosen, both literals are
false under the current assignment. Since $x^*$ satisfies the clause, at
least one of $l_1, l_2$ is true under $x^*$, so flipping that literal
decreases $d$ by~$1$. With probability at least $1/2$ the algorithm flips
a variable that decreases $d$. Therefore
\[
    \mathbb{E}[d(t+1) - d(t) \mid d(t) > 0] \leq -\frac{1}{2}.
\]
By a standard drift argument (see Chapter~5), the expected time to reach
$d = 0$ is at most $2n \cdot n = O(n^2)$.
\end{proof}

\begin{cpplisting}
\lstinputlisting[language=C++]{src/chapter6/main.cpp}
\end{cpplisting}

The key insight is that the random walk performs a biased random walk on
the integers $\{0, 1, \ldots, n\}$: from state $d > 0$ the expected drift
is negative, so the walk reaches~$0$ quickly. This
drift-condition argument generalizes beyond 2-SAT: whenever a random
process has a tendency to move toward a target state, hitting times are
often polynomial. We formalize this intuition through the theory of
Markov chains and conductance in subsequent sections.

%======================================================================
\section{Markov Chains}
\label{sec:markov-chains}
%======================================================================

\begin{definition}[Markov Chain]\label{def:markov-chain}
A \textbf{Markov chain} is a sequence of random variables
$X_0, X_1, X_2, \ldots$ taking values in a finite state space $\Omega$,
satisfying the \emph{Markov property}:
\[
    \Pr[X_{t+1} = j \mid X_t = i, X_{t-1} = i_{t-1}, \ldots, X_0 = i_0]
    = \Pr[X_{t+1} = j \mid X_t = i] = P_{ij}
\]
for all states $i, j \in \Omega$ and all $t \geq 0$. The matrix
$P = (P_{ij})_{i,j \in \Omega}$ is the \textbf{transition matrix}.
\end{definition}

A transition matrix is a row-stochastic matrix: $P_{ij} \geq 0$ for all
$i, j$ and $\sum_{j \in \Omega} P_{ij} = 1$ for all $i$. The entry
$P_{ij}$ gives the probability of moving from state~$i$ to state~$j$ in
one step. After $t$ steps, the probability of going from $i$ to $j$ is
the $(i,j)$-entry of $P^t$.

\begin{definition}[Stationary Distribution]\label{def:stationary}
A probability distribution $\pi$ on $\Omega$ is a \textbf{stationary
distribution} for the Markov chain with transition matrix $P$ if
\[
    \pi P = \pi,
\]
that is, $\sum_{i \in \Omega} \pi_i P_{ij} = \pi_j$ for all $j$. This
means that if the chain starts with distribution $\pi$, then at every
subsequent time step the distribution remains $\pi$.
\end{definition}

A Markov chain is \textbf{irreducible} if for every pair of states
$i, j$, there exists $t \geq 1$ with $(P^t)_{ij} > 0$: the chain can
reach any state from any other. It is \textbf{aperiodic} if
$\gcd\{t \geq 1 : (P^t)_{ii} > 0\} = 1$ for every state $i$. A chain
that is both irreducible and aperiodic is called \textbf{ergodic}.

\begin{theorem}[Convergence to Stationary Distribution]\label{thm:convergence}
Let $P$ be the transition matrix of an ergodic Markov chain on state
space $\Omega$ with $|\Omega| = n$. Then:
\begin{enumerate}[label=(\roman*)]
\item There exists a unique stationary distribution $\pi$ with $\pi_j > 0$
    for all $j$.
\item $\pi P = \pi$ (equilibrium equation).
\item For any initial distribution $\mu$,
    \[
        \lim_{t \to \infty} \mu P^t = \pi.
    \]
    That is, $P^t$ converges to the rank-one matrix $\mathbf{1}\pi$ as
    $t \to \infty$.
\end{enumerate}
\end{theorem}

The rate of convergence is governed by the second-largest eigenvalue of
$P$ (in absolute value), often called the \emph{spectral gap}
$\gamma = 1 - |\lambda_2|$. For reversible chains, a standard bound is
\[
    \|\mu P^t - \pi\|_{\mathrm{TV}} \leq \frac{1}{2} \sqrt{\frac{1}{\pi_{\min}}}\, |\lambda_2|^t.
\]

\begin{definition}[Detailed Balance]\label{def:detailed-balance}
The transition matrix $P$ with stationary distribution $\pi$ satisfies
the \textbf{detailed balance condition} if
\[
    \pi_i P_{ij} = \pi_j P_{ji} \quad \text{for all } i, j \in \Omega.
\]
A chain satisfying detailed balance is called \textbf{reversible}.
\end{definition}

Detailed balance is a sufficient condition for $\pi$ to be stationary:
summing both sides over $i$ gives $\pi_j = \sum_i \pi_i P_{ij}$.
Reversibility means that the chain looks the same when run forward or
backward in time. Not all stationary chains are reversible, but many
important chains arising from random walks on undirected graphs are.

\begin{example}[Random Walk on the Complete Graph]\label{ex:rw-complete}
Consider a random walk on the complete graph $K_n$. The transition matrix
is $P_{ij} = 1/n$ for all $i, j$ (including $P_{ii} = 1/n$). This chain
is ergodic. The uniform distribution $\pi = (1/n, \ldots, 1/n)$ is
stationary since $\sum_i (1/n)(1/n) = 1/n = \pi_j$. The chain satisfies
detailed balance because $\pi_i P_{ij} = 1/n^2 = \pi_j P_{ji}$. The
spectral gap is $\gamma = 1$, so the chain mixes in a single step.
\end{example}

\begin{example}[Lazy Random Walk]\label{ex:lazy-rw}
The \textbf{lazy random walk} on a graph $G$ stays at the current vertex
with probability $1/2$ and moves to a uniformly random neighbor with
probability $1/2$. The transition matrix is
$P = \frac{1}{2}I + \frac{1}{2}D^{-1}A$, where $A$ is the adjacency
matrix and $D$ is the diagonal degree matrix. Laziness ensures
aperiodicity without changing the stationary distribution.
\end{example}

%======================================================================
\section{Random Walks on Graphs}
\label{sec:rw-graphs}
%======================================================================

Let $G = (V, E)$ be an undirected graph with $n = |V|$ vertices and
$m = |E|$ edges. The simple random walk on $G$ is a Markov chain on $V$
where from vertex $v$, the walker moves to a uniformly random neighbor:
\[
    P_{uv} = \begin{cases} 1/d_u & \text{if } (u,v) \in E, \\ 0 &
    \text{otherwise}, \end{cases}
\]
where $d_u$ is the degree of $u$.

\begin{theorem}[Stationary Distribution for Graph Walks]\label{thm:rw-stationary}
The stationary distribution for the simple random walk on an undirected
graph $G$ with $m$ edges is
\[
    \pi_v = \frac{d_v}{2m}.
\]
Moreover, the chain satisfies detailed balance.
\end{theorem}

\begin{proof}
For an edge $(u,v) \in E$:
\[
    \pi_u P_{uv} = \frac{d_u}{2m} \cdot \frac{1}{d_u} = \frac{1}{2m}
    = \frac{d_v}{2m} \cdot \frac{1}{d_v} = \pi_v P_{vu}.
\]
Detailed balance holds, and $\pi$ is stationary.
\end{proof}

This theorem reveals that the random walk spends a fraction of time at
each vertex proportional to its degree. High-degree vertices are visited
more frequently. For $d$-regular graphs, the stationary distribution is
uniform.

\subsection{Mixing Time and Conductance}

The \textbf{total variation distance} between distributions $\mu$ and
$\nu$ on $\Omega$ is
\[
    \|\mu - \nu\|_{\mathrm{TV}} = \frac{1}{2} \sum_{x \in \Omega} |\mu(x) - \nu(x)|
    = \max_{S \subseteq \Omega} |\mu(S) - \nu(S)|.
\]
The \textbf{mixing time} is
\[
    t_{\mathrm{mix}}(\epsilon) = \min\left\{t : \|P^t(x, \cdot) -
    \pi\|_{\mathrm{TV}} \leq \epsilon \text{ for all } x \in V\right\},
\]
where $P^t(x,\cdot)$ is the distribution after $t$ steps starting from
$x$. A fundamental parameter controlling mixing is conductance.

\begin{definition}[Conductance]\label{def:conductance}
For a subset $S \subset V$ with $0 < \pi(S) \leq 1/2$, define
\[
    \Phi(S) = \frac{\text{Prob}(S \to \bar{S})}{\pi(S)}
    = \frac{\sum_{u \in S, v \in \bar{S}} \pi_u P_{uv}}{\pi(S)}.
\]
The \textbf{conductance} of $G$ is
\[
    \Phi(G) = \min_{S \subset V,\; 0 < \pi(S) \leq 1/2} \Phi(S).
\]
\end{definition}

Informally, conductance measures the bottleneck severity of the graph. A
graph with large conductance has no small bottlenecks and mixes quickly.
The conductance is also known as the \emph{Cheeger constant} of the graph
(in the discrete setting).

For a $d$-regular graph on $n$ vertices, the conductance is related to the
Cheeger inequality:
\[
    \frac{\gamma}{2} \leq \Phi(G) \leq \sqrt{2\gamma},
\]
where $\gamma = 1 - \lambda_2$ is the spectral gap.

%======================================================================
\section{Electrical Networks}
\label{sec:electrical-networks}
%======================================================================

There is a beautiful and deep correspondence between random walks on
graphs and electrical networks. Each edge $(u,v)$ of the graph is viewed
as a resistor of resistance $r_{uv} = 1$ ohm (for the unweighted case).
When a unit current is injected at a node $s$ and extracted at a node $t$,
the resulting voltages and currents encode information about the random
walk.

\begin{definition}[Effective Resistance]\label{def:eff-resistance}
The \textbf{effective resistance} $R_{\mathrm{eff}}(u,v)$ between nodes
$u$ and $v$ is the voltage difference between $u$ and $v$ when one unit
of current is injected at $u$ and extracted at $v$.
\end{definition}

Effective resistance satisfies the axioms of a metric: it is non-negative,
symmetric, $R_{\mathrm{eff}}(u,u) = 0$, and satisfies the triangle
inequality. It can be computed using Kirchhoff's laws or equivalently by
solving a system of linear equations derived from the graph Laplacian.

Let $L = D - A$ be the (combinatorial) Laplacian of $G$, where $D$ is the
diagonal degree matrix and $A$ is the adjacency matrix. The voltage
assignment $f$ satisfying $L f = I_{st}$ (where $I_{st}$ is the current
vector with $+1$ at $s$ and $-1$ at $t$) gives the effective resistance
via $R_{\mathrm{eff}}(s,t) = f(s) - f(t)$.

\begin{theorem}[Commute Time Identity]\label{thm:commute-time}
Let $G$ be an undirected graph with $m$ edges. The expected commute time
between nodes $u$ and $v$ is
\[
    \kappa(u,v) = H(u,v) + H(v,u) = 2m \cdot R_{\mathrm{eff}}(u,v),
\]
where $H(u,v)$ is the hitting time from $u$ to $v$.
\end{theorem}

This identity is remarkable: a probabilistic quantity (expected commute
time) equals a purely algebraic one ($2m \cdot R_{\mathrm{eff}}$) times a
constant. Several consequences follow:

\begin{corollary}\label{cor:hitting-upper}
For any $u, v \in V$,
\[
    H(u,v) \leq 2m \cdot R_{\mathrm{eff}}(u,v) \leq 2m \cdot n.
\]
\end{corollary}

\begin{corollary}[Maximal Commute Time]\label{cor:max-commute}
The maximum commute time satisfies $\max_{u,v} \kappa(u,v) \leq 2m \cdot n$,
and this bound is tight for the path graph.
\end{corollary}

The resistance metric also provides bounds on cover times and mixing times.
For instance, the commute time identity gives an alternative proof that the
random walk on a connected graph is recurrent in the sense that it visits
every vertex with probability~$1$.

%======================================================================
\section{Cover Times}
\label{sec:cover-times}
%======================================================================

The \textbf{cover time} $T_{\mathrm{cov}}(G)$ of a graph $G$ is the
expected number of steps for a random walk starting from the worst-case
vertex to visit every vertex at least once.

\begin{theorem}[Cover Time Bound]\label{thm:cover-time}
For any connected graph $G$ on $n$ vertices and $m$ edges,
\[
    T_{\mathrm{cov}}(G) \leq 2m(n - 1).
\]
\end{theorem}

\begin{proof}
By Theorem~\ref{thm:commute-time}, the commute time $\kappa(u,v) \leq
2m \cdot R_{\mathrm{eff}}(u,v) \leq 2m$. Order the vertices
$v_1, v_2, \ldots, v_n$ and consider the sequence of hitting times
$H(v_1, v_2), H(v_2, v_3), \ldots, H(v_{n-1}, v_n)$. The walk starting
at $v_1$ covers all vertices in time at most
\[
    \sum_{i=1}^{n-1} H(v_i, v_{i+1}) \leq (n-1) \cdot 2m.
\]
Taking the worst-case starting vertex gives the bound.
\end{proof}

\begin{theorem}[Matthews' Bound]\label{thm:matthews}
For any connected graph $G$,
\[
    T_{\mathrm{cov}}(G) \leq H_{\max} \cdot H_{n-1},
\]
where $H_{\max} = \max_{u,v} H(u,v)$ is the maximum hitting time and
$H_{n-1} = 1 + 1/2 + 1/3 + \cdots + 1/(n-1) = O(\ln n)$ is the
$(n-1)$-st harmonic number.
\end{theorem}

Matthews' bound is often tighter than the bound of
Theorem~\ref{thm:cover-time}. For many natural graph families, it gives
the correct asymptotic cover time up to constant factors.

\begin{example}[Cover Time of Specific Graphs]\label{ex:cover-times}
\begin{enumerate}[label=(\roman*)]
\item \textbf{Star graph} $K_{1,n-1}$: The cover time is
    $\Theta(n \ln n)$. The walk must visit all $n-1$ leaves, and each
    leaf is found in expected $O(1)$ steps from the center, but the
    walk must return to the center each time.
\item \textbf{Cycle} $C_n$: The cover time is $\Theta(n^2)$, matching
    the coupon collector for a one-dimensional random walk.
\item \textbf{Complete graph} $K_n$: The cover time is
    $\Theta(n \ln n)$, essentially the coupon collector problem.
\end{enumerate}
\end{example}

%======================================================================
\section{Graph Connectivity}
\label{sec:connectivity}
%======================================================================

Random walks provide a simple and efficient randomized algorithm for
testing whether a graph is connected. Given access to the graph via an
adjacency list oracle, the algorithm performs random walks and checks
whether all vertices are reached.

\begin{theorem}[Random Walk Connectivity]\label{thm:rw-connectivity}
Let $G$ be a connected, non-bipartite graph on $n$ vertices. A random
walk starting from any vertex visits every vertex within $O(n^3)$ steps
with high probability.
\end{theorem}

\begin{proof}
By Theorem~\ref{thm:cover-time} with $m \leq \binom{n}{2}$, the cover
time is at most $2m(n-1) \leq n^2 \cdot n = n^3$. A single walk of length
$O(n^3)$ therefore covers all vertices with constant probability. By
repeating $O(\log n)$ independent walks, the probability that some walk
fails to cover the graph is at most $n^{-c}$ for any constant~$c$.
\end{proof}

\begin{algorithm}[H]\label{alg:rw-connectivity}
\caption{Connectivity Testing via Random Walk}
\begin{algorithmic}[1]
\State Pick arbitrary start vertex $s$
\State Perform random walk from $s$ for $O(n^3)$ steps
\State Let $R$ = set of vertices visited
\If{$|R| = n$}
    \State \Return ``connected''
\Else
    \State \Return ``not connected''
\EndIf
\end{algorithmic}
\end{algorithm}

This algorithm uses $O(n^3)$ steps and $O(n)$ space, far less than the
$\Omega(n^2)$ space required to store the adjacency matrix. For sparse
graphs with $m = O(n)$, the running time improves to $O(n^2)$.

Bipartiteness can also be tested by random walks: a non-bipartite
connected graph has an odd cycle, and a random walk will detect this by
finding a vertex revisited at an odd distance from its first visit.

%======================================================================
\section{Expanders and Rapidly Mixing Random Walks}
\label{sec:expanders-bis}
%======================================================================

Expander graphs are sparse graphs with strong connectivity properties.
They play a central role in computer science, from error-correcting codes
to derandomization. The key property for our purposes is that random walks
on expanders mix to their stationary distribution in $O(\log n)$ steps.

\begin{definition}[Edge Expansion]\label{def:edge-expansion}
The \textbf{edge expansion} of a graph $G = (V, E)$ is
\[
    h(G) = \min_{S \subset V,\; |S| \leq n/2}
    \frac{|E(S, \bar{S})|}{|S|},
\]
where $E(S, \bar{S})$ is the set of edges crossing the cut.
\end{definition}

For a $d$-regular graph, the conductance $\Phi(G)$ equals
$h(G)/d$. Expander families have $\Phi(G) = \Omega(1)$, meaning the
conductance is bounded away from zero independently of~$n$.

\begin{theorem}[Expander Mixing Lemma]\label{thm:exp-mixing}
Let $G$ be a $d$-regular graph on $n$ vertices with eigenvalues
$d = \lambda_1 \geq \lambda_2 \geq \cdots \geq \lambda_n$. For any
$S, T \subseteq V$:
\[
    \left| |E(S,T)| - \frac{d \cdot |S| \cdot |T|}{n} \right|
    \leq \lambda_2 \sqrt{|S| \cdot |T|},
\]
where $E(S,T)$ counts edges between $S$ and $T$ (with edges inside
$S \cap T$ counted once).
\end{theorem}

The mixing lemma bounds the discrepancy between the actual number of
crossing edges and the expected number if edges were placed uniformly at
random. The bound is controlled by $\lambda_2$, the second-largest
eigenvalue: the smaller $\lambda_2$ (relative to $d$), the more
``random-like'' the edge distribution.

\begin{theorem}[Mixing Time of Expanders]\label{thm:exp-mixing-time}
Let $G$ be an expander graph with conductance $\Phi$ and spectral gap
$\gamma = 1 - |\lambda_2|/d$. For a random walk on $G$:
\[
    t_{\mathrm{mix}}(\epsilon) \leq \frac{1}{\gamma}
    \left(\ln \frac{1}{\pi_{\min}} + \ln \frac{1}{\epsilon}\right)
    = O\!\left(\frac{1}{\Phi^2} \log \frac{n}{\epsilon}\right).
\]
\end{theorem}

For constant-degree expanders, $\pi_{\min} = \Theta(1/n)$ and
$\Phi = \Theta(1)$, so the mixing time is $O(\log n)$. This is
exponentially faster than the $O(n^3)$ mixing time that suffices for
general connected graphs.

The existence of expander families is guaranteed by several
constructions. Random $d$-regular graphs are expanders with high
probability. Explicit constructions include Ramanujan graphs (which
achieve $|\lambda_2| \leq 2\sqrt{d-1}$, the Alon--Boppana bound) and
Lubotzky--Phillips--Sarnak graphs.

%======================================================================
\section{Probability Amplification by Random Walks on Expanders}
\label{sec:amplification}
%======================================================================

One of the most important applications of expander-based random walks is
\emph{probability amplification}: converting a weakly correct randomized
algorithm into one with arbitrarily small error probability, using very
few random bits.

\subsection{Almost-Uniform Sampling}

Consider a random walk on a constant-degree expander $G$ with $N$ vertices.
After $O(\log N)$ steps, the walk is within total variation distance
$\epsilon$ of the stationary distribution, which for a regular expander
is the uniform distribution. Therefore, the walk generates
approximately uniform random samples from $\{1, \ldots, N\}$ using only
$O(\log N)$ random bits per sample (to choose the initial vertex and the
transition at each step).

\begin{theorem}[Expander Random Walk Sampling]\label{thm:exp-sampling}
Let $G$ be an expander on $N$ vertices with conductance $\Phi > 0$ and
$\epsilon > 0$. A random walk of length $T = O(\frac{1}{\Phi^2} \log
\frac{1}{\epsilon})$ starting from any vertex produces a sample within
$\epsilon$ of uniform in total variation distance. Each sample uses
$O(\log N + T \log d)$ random bits, where $d$ is the degree.
\end{theorem}

\subsection{Error Reduction}

Suppose we have a randomized algorithm $A$ that, given a random string
$r \in \{0,1\}^k$, produces the correct answer with probability at least
$1/2 + \delta$ for some $\delta > 0$. We wish to boost this to
confidence $1 - \epsilon$.

The standard approach uses $O(k/\delta^2 \log(1/\epsilon))$ random bits
by taking the majority of $O(1/\delta^2 \log(1/\epsilon))$ independent
runs. Using an expander, we can reduce this to $O(k + \log(1/\epsilon) /
\delta^2)$ random bits as follows.

\begin{theorem}[Expander-Based Error Reduction]\label{thm:error-reduction}
Let $A$ be a randomized algorithm using $k$ random bits with success
probability $\geq 1/2 + \delta$. There exists a procedure that achieves
success probability $\geq 1 - \epsilon$ using
$O(k + \frac{1}{\delta^2} \log \frac{1}{\epsilon})$ random bits.
\end{theorem}

\begin{proof}
Let $G$ be an expander on $2^k$ vertices with degree $d$ and second
eigenvalue $|\lambda_2| \leq d/4$ (such expanders exist for constant
$d$). Enumerate the vertices of $G$ with binary strings of length $k$.
Start a random walk at a uniformly random vertex $v_0 \in \{0,1\}^k$.
After $T$ steps, we have visited vertices $v_0, v_1, \ldots, v_T$, and we
run $A$ on each, taking the majority answer.

Since $G$ is an expander, the random walk produces $T$ vertices that are
approximately independent from the perspective of $A$'s success
criterion. More precisely, using the expander mixing lemma, the
correlation between the events ``$A$ succeeds at $v_i$'' and ``$A$
succeeds at $v_j$'' decays exponentially with the graph distance between
$v_i$ and $v_j$ in the walk.

By standard results on expander walks, for any predicate $B$ on
$\{0,1\}^k$ with $\Pr_{r \leftarrow \{0,1\}^k}[B(r)] \geq 1/2 + \delta$,
the probability that the majority of $T = O(1/\delta^2)$ consecutive
vertices on the walk satisfy $B$ is at least $1 - e^{-\Omega(\delta^2
T)}$. Setting $T = c/\delta^2 \cdot \ln(1/\epsilon)$ gives failure
probability at most $\epsilon$, using only $O(k + T \log d)$ random bits.
\end{proof}

The key advantage over independent sampling is that the expander walk
uses $O(k + T \log d)$ random bits total, whereas independent sampling
requires $O(kT)$ bits. The expander walk ``reuses'' its random bits
because each step of the walk requires only $\log d$ new bits to choose
which edge to follow.

\subsection{Applications}

Expander-based probability amplification has numerous applications:

\begin{enumerate}[label=(\roman*)]
\item \textbf{BPP $\subseteq$ P/poly}: Expander walks reduce the random
    seed length, enabling the construction of polynomial-size
    advice strings that make the algorithm deterministic.

\item \textbf{Pseudorandom generators}: Expander-based constructions
    yield PRGs that fool any family of small circuits, connecting
    to derandomization of BPP.

\item \textbf{Sampling and counting}: In the context of Markov chain
    Monte Carlo (MCMC), rapidly mixing chains on expanders provide
    efficient approximate samplers for complex combinatorial
    distributions (Jerrum--Sinclair).
\end{enumerate}

%======================================================================
\addcontentsline{toc}{section}{Notes}
\section*{Notes}
%======================================================================

The theory of Markov chains and random walks on graphs has rich origins
in both probability theory and combinatorics. The foundational survey by
Aldous~\cite{Aldous1983} established the hitting time, cover time, and
commute time framework and identified many key open problems. The
comprehensive treatment by Lov\'{a}sz~\cite{Lovasz1993} systematically
developed the spectral theory of random walks on graphs, including
detailed proofs of the stationary distribution theorem, the Cheeger
inequality, and mixing time bounds.

The connection between random walks and electrical networks, developed by
Lyons and Peres~\cite{LyonsPeres2016} (building on earlier work by
Doob, Kahane, and others), is presented in their textbook \emph{Probability
on Trees and Networks}. Theorem~\ref{thm:commute-time} appears there as
the ``commute time identity.''

Matthews' bound (Theorem~\ref{thm:matthews}) was proved by
Matthews~\cite{Matthews1988} and independently by Aldous. The tight
cover time results for specific graph families were established in a
series of papers by Aldous~\cite{Aldous1989, Aldous1991}.

The use of random walks for testing graph connectivity dates to
Aleliunas~\emph{et al.}~\cite{AleliunasEtAl1979}, who proved the
$O(n^3)$ bound (Theorem~\ref{thm:rw-connectivity}). This was one of the
first applications of random walks in algorithm design.

The theory of expander graphs was initiated by Pinsker~\cite{Pinsker1973}
and Margulis~\cite{Margulis1973}, who proved the existence of expander
families. Explicit constructions were given by
Margulis~\cite{Margulis1988}, Gabai and Zemor, and
Lubotzky--Phillips--Sarnak~\cite{LubotzkyEtAl1988}. The Ramanujan
bound $|\lambda_2| \leq 2\sqrt{d-1}$ is achieved by Ramanujan graphs.
The expander mixing lemma (Theorem~\ref{thm:exp-mixing}) is due to
Alon and Chung~\cite{AlonChung1988}.

The use of expander walks for probability amplification was developed by
Impagliazzo, Zuckerman, and Wigderson~\cite{ImpagliazzoEtAl1995}.
Jerrum and Sinclair~\cite{JerrumSinclair1996} pioneered the use of
rapidly mixing Markov chains for approximate counting and sampling, a
technique with broad applications in statistical physics, combinatorics,
and machine learning.

The random walk approach to 2-SAT (Section~\ref{sec:2sat-example}) was
introduced by Papadimitriou~\cite{Papadimitriou1991} and analyzed by
Sch\"oning~\cite{Schoninger1993} in the context of $k$-SAT algorithms.

%======================================================================
\addcontentsline{toc}{section}{Problems}
\section*{Problems}
%======================================================================

\begin{problem}
Prove that for a lazy random walk on any connected undirected graph, the
chain is aperiodic. Show that the stationary distribution is the same as
for the non-lazy walk: $\pi_v = d_v / 2m$.
\end{problem}

\begin{problem}
Compute the stationary distribution and the spectral gap for the simple
random walk on the cycle $C_n$. What is the mixing time $t_{\mathrm{mix}}(1/4)$?
\end{problem}

\begin{problem}
Let $G$ be a $d$-regular graph on $n$ vertices with eigenvalues
$d = \lambda_1 \geq \lambda_2 \geq \cdots \geq \lambda_n$. Prove that the
conductance satisfies $\Phi(G) \leq \sqrt{2(1 - \lambda_2/d)}$. This is
the Cheeger inequality for regular graphs.
\end{problem}

\begin{problem}
Verify the commute time identity (Theorem~\ref{thm:commute-time}) for the
path graph $P_n$ with vertices $1, 2, \ldots, n$ and edges $(i, i+1)$.
Compute $H(1, n)$, $H(n, 1)$, and $R_{\mathrm{eff}}(1, n)$ explicitly.
\end{problem}

\begin{problem}
Prove that the cover time of the complete graph $K_n$ is $n H_{n-1} =
n \ln n + O(n)$. (Hint: model the cover time as a coupon collector
problem.)
\end{problem}

\begin{problem}
Let $G$ be an $n$-vertex graph with maximum degree $\Delta$. Prove that the
cover time satisfies $T_{\mathrm{cov}}(G) \geq n(n-1)/(2\Delta)$.
Give an example where this bound is tight up to a constant factor.
\end{problem}

\begin{problem}
Prove that for a random walk on a connected non-bipartite graph, the
probability of being at vertex $v$ after $t$ steps converges to $\pi_v$
at rate $O(|\lambda_2|^t)$, where $\lambda_2$ is the second-largest
eigenvalue of the transition matrix in absolute value.
\end{problem}

\begin{problem}
Show that the expander mixing lemma implies the following: for a
$d$-regular expander $G$ on $n$ vertices with $|\lambda_2| \leq d/4$, any
set $S$ with $|S| \leq n/2$ satisfies $|E(S, \bar{S})| \geq
d|S|/8$.
\end{problem}

\begin{problem}
\textbf{(Biased Random Walk.)} Consider a random walk on the path
$\{0, 1, \ldots, N\}$ that moves right with probability $p$ and left
with probability $q = 1 - p$, with absorbing barriers at $0$ and $N$.
For $p \neq q$, compute the hitting time $H(0, N)$ exactly. Compare with
the unbiased case $p = q = 1/2$.
\end{problem}

\begin{problem}
Let $G$ be a $d$-regular expander on $N$ vertices. Design a randomized
algorithm that, given an oracle for a function $f : \{0,1\}^k \to
\{0,1\}$ (where $N = 2^k$), estimates the bias $\Pr[f = 1]$ to within
additive error $\epsilon$ using $O(k/\epsilon^2)$ queries to $f$.
(Hint: use a random walk on $G$ to generate correlated samples and apply
the expander mixing lemma to bound the variance of the estimator.)
\end{problem}

\end{document}
